\section{Discussion and open questions}
\label{sec:discussion}

The exact enumerators change the interpretation of EA and EC parameters.
There is no single mixing bottleneck.  Dense messages already behave like the
input to a random invertible map and attain the GV exponent.  Sparse messages
determine how many expander edges or convolution steps are needed before that
regime becomes available.

\paragraph{What convolution buys.}
The wrapped convolution improves the sparse endpoint from
$I_\delta=1-2\sqrt{\delta(1-\delta)}$ for the accumulator to
$I^{\mathrm w}_{m,\delta}$ in \eqref{eq:ec-sparse-rate}.  The latter tends to
$1-2\delta$ as memory grows.  It does not improve the dense exponent beyond
the random-linear-code value.  This explains the finite frontier: a few
memory words can replace several expander edges, but large memory eventually
has diminishing returns.

\paragraph{What regularity buys.}
Left regularity forces every message coordinate to contribute once per
region.  It reduces the variance in edge placement that dominates low-support
terms.  Two-sided regularity also eliminates empty right coordinates at full
support.  This property is essential to the large-field certificates.  Over
the binary field, odd right degree also prevents the all-one message from
entering the expander kernel.  Two-sided regularity substantially improves
the binary EC profile, but it does not improve the weight-two EA bottleneck.

The certified frontier is not exhaustive.  Remaining odd right degrees may
yield additional tradeoffs.  Measured encoding costs should eventually
replace the proxy $d_L+m$.  Near-regular right degrees are another possible
way to avoid the binary parity obstruction, but they require a new point-mass
bound.

\paragraph{Binary choices versus field values.}
The binary profiles target Silent OT, whose choices are bits.  If the same
matrix is viewed over $\F_{2^{128}}$, Lemma~\ref{lem:binary-scalar-extension}
still preserves its Hamming distance, but the map splits into 128 independent
binary maps.  Silent VOLE instead needs mixing between these components.  Our
field-valued profile supplies that mixing through ambient-field edge labels
and full-field convolution coefficients.

\paragraph{Field size.}
Random nonzero edge labels are not cosmetic.  They make the output law depend
only on message support and enable the projective trace bounds.  They also add
one multiplication per edge.  Finite-field EC additionally needs full-field
feedback coefficients for the reduced-state argument.  A more restricted tap
distribution may be faster, but it requires a state description that retains
more than the zero pattern.

The finite failure exponents can decrease with the field order when the target
is a fixed fraction of the field-relative GV distance.  Consider field order
$p=2^s$, where $s$ is the bit length of the field order.  At rate one half,
\[
 \delta_{2^s,\mathrm{GV}}(1/2)=\frac12-\frac1s+O(s^{-2}).
\]
Near $s=128$, increasing $s$ by one raises the GV-relative cutoff by about
$126$ output coordinates at our length.  The low-equation
structural branch receives no compensating field factor after its projective
cap saturates.  At a fixed absolute cutoff, increasing $s$ does not weaken
this branch; the field-suppressed branch only becomes smaller.

Table~\ref{tab:field-size-diagnostic} illustrates the distinction for the
degree-$24/12$, memory-$5$ certificate.  It freezes the $\F_{2^{128}}$ support
partition and positive markers.  It then uses either the field-relative
cutoff from Theorem~\ref{thm:field-size-ec-certificates} or the fixed
$\F_{2^{128}}$ cutoff.  The frozen exponent is the negative base-two
logarithm of this floating-point bound.  The certified column reports
outward-rounded interval bounds where available.  The smallest stored field
marker is $2^{-127}$, so every reused marker remains valid for the displayed
orders.

\begin{table}[ht]
\centering
\caption{Field-size profile around 128-bit fields for degrees $24/12$ and
memory $5$.  A dash denotes a frozen-marker diagnostic without a separate
interval certificate.  The last column keeps the $\F_{2^{128}}$ cutoff
fixed.}
\label{tab:field-size-diagnostic}
\begin{tabular}{@{}rrrrr@{}}
\toprule
$s$ & GV-relative cutoff & frozen bits & certified bits & fixed-cutoff bits \\
\midrule
$127$ & $1{,}015{,}699$ & $45.84$ & $45.84$ & $45.03$ \\
$128$ & $1{,}015{,}826$ & $45.03$ & $45.03$ & $45.03$ \\
$129$ & $1{,}015{,}950$ & $44.24$ & $44.24$ & $45.03$ \\
$130$ & $1{,}016{,}073$ & $43.44$ & --- & $45.03$ \\
$132$ & $1{,}016{,}314$ & $41.84$ & --- & $45.03$ \\
$136$ & $1{,}016{,}774$ & $38.69$ & $38.78$ & $45.03$ \\
$140$ & $1{,}017{,}207$ & $35.63$ & --- & $45.03$ \\
$144$ & $1{,}017{,}616$ & $32.67$ & --- & $45.03$ \\
\bottomrule
\end{tabular}
\end{table}

The certificates at $s=127,128,129$ confirm a local loss of about $0.80$
failure bits per additional field bit.  At $s=136$, local retuning improves
the frozen exponent from $38.6903$ to $38.7836$ bits.  The gain is only
$0.0933$ bits.

The structural branch controls this regime.  At $s=136$, the largest exact
band covers message supports $129$ through $160$.  Its frozen structural and
field-suppressed exponents are $-41.3863$ and $-162.9032$, respectively.
Thus additional field suppression cannot compensate for the larger cutoff.
At the fixed $\F_{2^{128}}$ cutoff $L=1{,}015{,}826$, every displayed field
order rounds to the same $45.03$-bit diagnostic exponent.  The decline comes
from moving the GV-relative target, not from weaker mixing at the original
distance.

This comparison does not include the certified $\F_{2^{127}-1}$ profile.
That profile uses different degree--memory parameters and the full floored
GV cutoff.

Finding a faster feedback distribution with an exact small state space
remains open.  A useful candidate must retain enough symmetry for a rigorous
trace bound and use fewer field operations than full-field random taps.

\paragraph{Implementation boundary.}
The proved finite-field ensemble assigns an independent nonzero label to every
expander edge.  These labels act before the incident edges are summed.  They
therefore control cancellations at an expander coordinate.

Our reference implementation follows this structure.  It materializes an
independent balanced assignment for each region and gives every edge an
independent nonzero label.  It also samples every convolution tap from the
full field.  A public seed expands all three sources of randomness.

The binary implementation materializes the two-sided regular expander and
the wrapped convolution analyzed in Section~\ref{sec:regular}.  It is the
Silent OT profile, not the Silent VOLE profile.

Our streaming implementation studies a different construction.  Over
odd-characteristic fields it uses a signed structured expander.  Over
characteristic two it uses compact full-index permutations.  Both paths
regenerate full-field convolution coefficients from a public seed and use
field multiplications in the recurrence.  These multiplications provide the
component mixing needed by field-valued VOLE.  The expander labels are not
independent nonzero field elements, so we treat the streaming construction as
heuristic rather than as a sample from the proved ensemble.

Our experiments also include unit edge coefficients as an ablation.  They do
not provide evidence of a defect at the tested sizes, but this observation is
not a distance proof.  Proving useful distance for structured or unit edge
coefficients remains a separate problem.

\paragraph{Cryptographic interpretation.}
A distance certificate bounds the probability that code generation produces
a low-weight nonzero codeword.  It does not by itself establish the security
of a PCG or a silent protocol.  Such an application must connect its attack
model and noise distribution to the certified distance event.  It must also
account for any public structure not sampled by our ensemble.

For Silent VOLE over $\F_{2^{128}}$, the finite-field certificate supplies a
distance premise for the independently labeled reference ensemble.  A
complete protocol theorem must still connect that premise to the chosen noise
distribution and attack model.  It must separately justify the structured
streaming heuristic used by the implementation.

The following composition rule records exactly how the sampling term enters a
protocol theorem.

\begin{proposition}[Adding code-sampling failure]
\label{prop:application-composition}
Let $\G$ be a sampled generator, define the set
$\mathcal G_{\mathrm{good}}:=\{G:d_{\min}(G)>L\}$, and suppose
$\Pr[\G\notin\mathcal G_{\mathrm{good}}]\leq U$.  Let $\mathsf{Fail}$ be an event defined on
the joint randomness of code generation and a protocol.  If
\begin{equation}
 \Pr[\mathsf{Fail}\mid\G=G]\leq\varepsilon
 \qquad\text{for every }G\in\mathcal G_{\mathrm{good}},
\end{equation}
then $\Pr[\mathsf{Fail}]\leq\varepsilon+U$.
\end{proposition}

\begin{proof}
Condition on whether $\G\in\mathcal G_{\mathrm{good}}$.  The good branch
contributes at most $\varepsilon$, and the bad branch contributes at most
$U$.
\end{proof}

The same additive term applies to a distinguishing advantage by replacing
the bad-code branch with the trivial upper bound one.  A Silent OT or Silent
VOLE analysis must still prove the conditional bound $\varepsilon$ for its
noise distribution and attack model.  This paper supplies only $U$.

\paragraph{Decoding and explicit constructions.}
The proofs are ensemble-existence and sampling statements.  They do not give
an efficient decoding algorithm.  They also do not produce a deterministic
family with the same distance.  Both questions are orthogonal to the present
first-moment method.

\paragraph{Certificate maintenance.}
The current artifacts use several specialized verifiers.  A shared library
for interval partitions, Arb matrices, and positive-marker bounds would make
the trusted implementation smaller.  An independent verifier in a second
language would further reduce implementation risk.

The checked-in manifest maps each claimed numerical result to its verifier,
inputs, SHA-256 digests, arithmetic precision, and expected failure exponent.
The manifest driver checks these records and can execute all rigorous entries
serially.  The specialized verifiers remain part of the trusted base.  A
second implementation would reduce that trust further.

\paragraph{Conclusion.}
Exact input--output enumerators give substantially stronger guarantees for
the original EA and EC ensembles.  In the binary case, region-stratified
regularity controls sparse messages and gives a concrete rate-one-half
degree--memory frontier near the GV cutoff.  In fields near 128 bits, random
edge labels, projective grouping, and two-sided regularity control the
additional union over message values.  The finite-field analysis gives a
separately certified field-mixing profile over $\F_{2^{127}-1}$.  It also
gives one fixed degree--memory profile over field orders $2^{127}$ through
$2^{136}$ at a concrete cryptographic-scale length.
